\chapter{Padronização e estilo}

\section{Figuras, tabelas e gráficos}

\begin{itemize}
\item tabela = quadro;
\item figuras devem ter chamada no texto (de acordo com o projeto gráfico);
\item tabelas e figuras devem ter legendas (de acordo com o projeto gráfico);
\item desenho e diagramação de tabelas e figuras obedecem ao projeto gráfico;
\item tabelas extensas no formato paisagem;
\item créditos de autoria são obrigatórios e devem ser postos no entorno da figura/tabela ou na mesma página, quando se referir a múltiplas figuras/tabelas;
\item licença (quando for o caso) deve constar junto do crédito.
\end{itemize}

\section{Titulação}

(Por produto)

\section{Paginação}

\section{Destaque}

\subsection{Negrito}

\subsection{Citações}

\subsubsection{Aspas duplas}

Eventos devem ser destacados com aspas duplas, para que possam ser postados em redes sociais.

\subsubsection{Aspas simples}

Para evitar o uso do itálico (v. \ref{acessibilidade}), palavras listadas como estrangeiras no \emph{Vocabulário ortográfico da língua portuguesa} (Volp) devem vir entre aspas simples.

\subsubsection{Recuos marginais}

\subsection{Itálico}

\section{Ortografia}

\subsection{Maiúsculas e minúsculas}

Como regra geral, usam-se maiúsculas e minúsculos para nomes próprios de pessoas, documentos e instituições:

\begin{itemize}
\item Conselho Regional de Psicologia;
\item Prefeitura Municipal de São Paulo.
\end{itemize}

Em todos os outros casos, usam-se apenas minúsculas:

\begin{itemize}
    \item carteira de identidade profissional (comum a todas as profissões) --- confirmar;
\item pessoa física, pessoa jurídica (quando não forem parte do nome de um documento --- Cadastro de Pessoa Física, por exemplo)
\end{itemize}

\subsubsection{Exceções}

\emph{Conselho}, quando se tratar do Conselho Regional de Psicologia de São Paulo.

\paragraph{Áreas de estudo}

Em caixa alta e baixa quando se referir a grandes áreas (Psicologia, Educação, Medicina) e em caixa-baixa quando a subáreas (psicologia clínica); em caixa alta e baixa quando se tratar de bandeiras do CRP (Direitos Humanos, Políticas Públicas --- ?). 

\paragraph{Títulos acadêmicos}

Usa-se o título apenas para os graus de mestra/e e doutora/doutor. Para especialização (pós-graduação `lato sensu') e pós-doutorado, usa-se ``com especialização em'' ou ``com pós-doutorado em''.

\begin{itemize}
\item Títulos acadêmicos devem vir sempre acompanhados da área de estudo e da instituição que emite o diploma (p. ex., "mestra em Psicologia Escolar e do Desenvolvimento Humano pela Universidade de São Paulo");
\item o título de \emph{especialista} é reservado para psicólogas/os que obtiveram registro de especialista junto ao CFP.
\end{itemize}

\subsection{Sinais gráficos}

\subsubsection{Barra}

Atendimento/Pessoa jurídica (seguindo a lógica das pastas no computador).

Não deve haver espaço entre a barra (/) e as letras que a circundam..

\subsubsection{Hífen}

\paragraph{Cargos}

O hífen deve ser usado quando o termo composto define uma posição hierárquica: vice-presidente, coordenador-adjunto etc.

Não se usa hífen quando se trata de funções simultâneas, como conselheira presidente.

\subsection{Parêntese}

\subsection{Grafia oficial das comissões e produtos de comunicação do CRP~SP}

\subsubsection{Conselho Regional de Psicologia da 6ª Região}

O nome CRP~SP \emph{sempre} deve ser grafado com um espaço incondicional (Alt+0160) entre a sigla do Conselho e a do estado. Isso evita a translineação, quer dizer, que o nome seja dividido pelos programas de edição de texto (CRP numa linha e SP em outra).

O número indicativo da região (6ª Região) deve ser feito entre parênteses (trata-se de informação complementar).

\subsubsection{Comissões}

\begin{itemize}
    \item[\textbf{Caci}]{Comissão de auditoria e controle interno}
	\item[\textbf{CamCoe}]{Câmara de Mediação da COE}
	\item[\textbf{Carpe}]{Comissão de análise de registro de Psicólogo/a especialista}
	\item[\textbf{CDH}]{Comissão Especial de Direitos Humanos}
	\item[\textbf{CHM}]{Comissão Especial de História e Memória da Psicologia}
	\item[\textbf{CL}]{Comissão Permanente de Licitação}
	\item[\textbf{COE}]{Comissão de Ética}
	\item[\textbf{COF}]{Comissão de Orientação e Fiscalização}
	\item[\textbf{ComCom}]{Comissão de Comunicação}
	\item[\textbf{CPLC}]{Comissão Especial Processo Legislativo e Concursos}
	\item[\textbf{Crepop}]{Centro de referências técnicas em Psicologia e Políticas Públicas}
	\item[\textbf{Cred}]{Comissão Especial Riscos, Emergências e Desastres}
	\item[\textbf{Crer}]{Comissão Especial Relações Étnico-Raciais}
\end{itemize}

\subsubsection{Produtos de comunicação}

\section{Itens, alíneas, enumerações}

\subsection{listas e enumerações}

\begin{itemize}
    \item Podem ser utilizadas quando houver \href{https://www.w3.org/TR/wcag-3.0/text-and-wording}{três ou mais itens relacionados}
    \item    documentos, informações reunidas etc.
\end{itemize}

\subsubsection{Listas não numeradas}

Utilizar `bullets'.

\subsubsection{Listas numeradas}
 
Utilizadas para passo a passo, processos que envolvem diversas etapas etc.

\section{Numerais}

\subsection{Cardinais e ordinais}

\subsubsection{Ordinais}

feminino (1ª): Alt + 166
masculino (1º): Alt + 167

\subsection{Sinais matemáticos}

\subsubsection{Oposição}

× (multiplicação ou oposição — m.q. *vs.*) = Alt + 0215

\subsection{Numerais em leis, portarias, resoluções etc.}

\subsection{Numerais com unidades de medida}

\subsection{Chamadas de nota}

devem ser posicionadas após a pontuação

\section{Siglas e abreviações}

\subsection{Abreviações}

Nomes abreviados não trazem espaço entre o ponto e a letra abreviada seguinte (S.J. dos Campos).

\subsection{Siglas}

Siglas não devem ter o plural flexionado (as ONG, os CRP etc.).

\section{Gramática}

\subsection{Plural}

Atenção a coletivos, metonímias e termos abstratos que não devem ser flexionados no plural: \emph{legislação} (e não ``legislações''), \emph{material} (evitar ``materiais'') e \emph{violência} (e não ``violências'').

\section{Acessibilidade}\label{acessibilidade}

\begin{enumerate}
\item \href{https://axesslab.com/fonts-dont-matter/}{incluir diversos títulos e subtítulos} e listas não numeradas, para que se evitem grandes blocos de texto (preferir parágrafos curtos);
\item complementar o texto com infográficos, vídeos, imagens, áudio;
\item evitar o uso de itálico;
\item atentar para a extensão das linhas, que deve ser de 50 a 75 caracteres (ou entre 40 e 50 para texto em duas colunas)
\begin{itemize}
\item altura da linha de 150\% do tamanho da fonte;
\item fontes de 18 pixels no mínimo (para web);
\end{itemize}
\item utilizar fontes populares;
\item considerar a inclusão de glossários e listas de siglas nos documentos;
\item utilizar sinais gráficos corretos (reticências etc.);
\item textos alternativos para imagens (exceto quando a imagem for decorativa/ilustrativa e não trouxer informações que não se encontram no texto) e de conteúdos gerados externamente ao InDesign (no Illustrator, por exemplo);
\item incluir títulos em tabelas;
\item não confiar apenas nas cores para elaborar legendas (em gráficos, por exemplo);
\item utilizar \href{https://www.gov.br/governodigital/pt-br/acessibilidade-e-usuario/acessibilidade-digital/eMAGv31.pdf}{âncoras de conteúdo} (minissumários no início de cada seção autônoma);
\item\href{https://www.gov.br/governodigital/pt-br/acessibilidade-e-usuario/acessibilidade-digital/eMAGv31.pdf}{não utilizar tabelas para diagramação};
\item\href{https://www.nngroup.com/articles/web-writing-show-numbers-as-numerals/}{números podem ser escritos por extenso em material impresso/PDF, mas [devem ser grafados em algarismos arábicos na web};
\item\href{https://www.gov.br/governodigital/pt-br/acessibilidade-e-usuario/acessibilidade-digital/eMAGv31.pdf}{identificar o idioma principal do documento e informar quando houver mudança}.
\end{enumerate}

\section{Toponímia}

\subsection{Logradouros}

Seguindo o princípio de acordo com o qual usa-se caixa alta e baixa em nomes próprios (``Programa de Pós-Graduação em Psicologia''), podemos adotar "rua" quando não houver nome que a caracterize (``as ruas de São Paulo'', ``na rua do CRP'') e ``Rua Teodoro Sampaio'' quando for seguida do nome. Ou não.

\subsubsection{Estados}

Seguindo o princípio adotado nas regras de gendramento segundo o qual a barra indica alternância, indicam-se os nomes de estados entre parênteses:

Lorena (SP)

Obs.: nas normas de referências bibliográficas da ABNT, cidade e estado são separados por vírgulas:

Lorena, SP



